\documentclass[11pt,a4paper]{article}

\usepackage[margin=1in]{geometry}
\usepackage{amsmath,amssymb,amsthm}
\usepackage{mathtools}
\usepackage{bm}
\usepackage{braket}
\usepackage{hyperref}
\usepackage{xcolor}
\usepackage{enumitem}

\hypersetup{
  colorlinks=true,
  linkcolor=blue!60!black,
  urlcolor=blue!60!black,
  citecolor=blue!60!black
}

\newcommand{\rr}{\mathbf{r}}
\newcommand{\kk}{\mathbf{k}}
\newcommand{\GG}{\mathbf{G}}
\newcommand{\RR}{\mathbf{R}}
\newcommand{\TT}{\mathbf{T}}
\newcommand{\tautau}{\boldsymbol{\tau}}
\newcommand{\tvec}{\mathbf{t}}
\newcommand{\kibz}{\kk_{\text{ibz}}}
\newcommand{\Uhat}{\hat{U}}

\title{Symmetry Relations for PAW Pseudo-Orbitals \\
       at Symmetry-Related $k$-points}
\author{Detailed derivation notes}
\date{}

\begin{document}
\maketitle

\tableofcontents
\bigskip\hrule\bigskip

\section{Preliminaries and notation}

We work with a periodic crystal whose Bravais lattice is generated by primitive
vectors $\{\mathbf{a}_1,\mathbf{a}_2,\mathbf{a}_3\}$, with reciprocal lattice
spanned by $\{\mathbf{b}_1,\mathbf{b}_2,\mathbf{b}_3\}$. Atoms in the home unit
cell sit at basis positions $\tautau_a$, $a = 1,\dots,N_{\text{at}}$. Lattice
vectors are denoted $\RR$, reciprocal lattice vectors $\GG$, and crystal
momenta $\kk$ in the first Brillouin zone (BZ).

A space-group operation is written
\begin{equation}
g = \{R \mid \tvec\}, \qquad g\,\rr = R\,\rr + \tvec,
\end{equation}
where $R \in O(3)$ is a point-group rotation (or improper rotation) and $\tvec$
is a possibly non-zero fractional translation (nonsymmorphic case). The action
on a wavefunction is the standard scalar-field action,
\begin{equation}
\bigl[\Uhat(g)\,\psi\bigr](\rr) \;=\; \psi(g^{-1}\rr) \;=\; \psi\bigl(R^{-1}(\rr - \tvec)\bigr),
\end{equation}
which is unitary on $L^2(\mathbb{R}^3)$.

We treat the spin-degenerate, non-relativistic case throughout; comments on
spin--orbit coupling and time reversal appear at the end.

\subsection{The PAW transformation}

The PAW formalism\footnote{P.~E.~Bl\"ochl, \emph{Phys.~Rev.~B} \textbf{50},
17953 (1994).} relates an all-electron (AE) Bloch orbital
$\ket{\psi_{n\kk}}$ to a smooth pseudo-orbital $\ket{\tilde\psi_{n\kk}}$ via
a linear operator $\mathcal{T} = \mathbf{1} + \sum_a \mathcal{T}_a$, where each
$\mathcal{T}_a$ is supported in an augmentation region around atom $a$.
Concretely,
\begin{equation}
\ket{\psi_{n\kk}}
\;=\; \ket{\tilde\psi_{n\kk}}
+ \sum_{\RR}\sum_{a,i}
   \Bigl(\ket{\phi_i^{a,\RR}} - \ket{\tilde\phi_i^{a,\RR}}\Bigr)
   \langle \tilde p_i^{a,\RR} \mid \tilde\psi_{n\kk}\rangle,
\label{eq:paw-real}
\end{equation}
where the channel index $i = (n_i,\ell_i,m_i)$ collects the radial and angular
quantum numbers, and the localized objects at site $(a,\RR)$ are
\begin{align}
\phi_i^{a,\RR}(\rr)
   &= R_{n_i\ell_i}^{a}\!\bigl(|\rr - \tautau_a - \RR|\bigr)\,
      Y_{\ell_i m_i}\!\bigl(\widehat{\rr - \tautau_a - \RR}\bigr), \\
\tilde\phi_i^{a,\RR}(\rr)
   &= \tilde R_{n_i\ell_i}^{a}\!\bigl(|\rr - \tautau_a - \RR|\bigr)\,
      Y_{\ell_i m_i}\!\bigl(\widehat{\rr - \tautau_a - \RR}\bigr), \\
\tilde p_i^{a,\RR}(\rr)
   &= \tilde p_{n_i\ell_i}^{a}\!\bigl(|\rr - \tautau_a - \RR|\bigr)\,
      Y_{\ell_i m_i}\!\bigl(\widehat{\rr - \tautau_a - \RR}\bigr).
\end{align}
The radial functions $R^a$, $\tilde R^a$, $\tilde p^a$ depend only on the atomic
species; the projectors $\tilde p_i^{a,\RR}$ are dual to the pseudo partial
waves, $\langle \tilde p_i^{a,\RR}\mid \tilde\phi_j^{a',\RR'}\rangle =
\delta_{aa'}\delta_{\RR\RR'}\delta_{ij}$.

\subsection{Bloch sums and $k$-dependent projectors}

The localized projector $\tilde p_i^{a,\RR}$ has \emph{no} $\kk$-dependence;
it is an atomic-like real-space function attached to a particular cell. The
object that pairs naturally with a Bloch state at $\kk$ is the Bloch-summed
projector
\begin{equation}
\ket{\tilde p_i^{a,\kk}}
\;=\; \sum_{\RR} e^{i\kk\cdot(\tautau_a + \RR)}\,\ket{\tilde p_i^{a,\RR}},
\label{eq:bloch-projector}
\end{equation}
and similarly for $\phi_i^{a,\kk}$ and $\tilde\phi_i^{a,\kk}$. The
phase convention chosen here puts the basis-position phase
$e^{i\kk\cdot\tautau_a}$ into the projector itself (the so-called
\emph{periodic} or \emph{convention II} gauge); a common alternative is to
keep only the lattice phase $e^{i\kk\cdot\RR}$ and let the basis phase live
elsewhere. We will be explicit about this convention whenever it matters.

With \eqref{eq:bloch-projector}, the PAW expansion \eqref{eq:paw-real} reduces
to a \emph{home-cell} sum:
\begin{equation}
\boxed{\;\ket{\psi_{n\kk}}
   \;=\; \ket{\tilde\psi_{n\kk}}
   + \sum_{a,i}\Bigl(\ket{\phi_i^{a,\kk}} - \ket{\tilde\phi_i^{a,\kk}}\Bigr)\,
                 c_i^{a}(n,\kk),\;}
\label{eq:paw-bloch}
\end{equation}
with projector coefficients
\begin{equation}
c_i^a(n,\kk) \;\equiv\; \langle \tilde p_i^{a,\kk} \mid \tilde\psi_{n\kk}\rangle.
\label{eq:projector-coef}
\end{equation}

The $\kk$-label on $\ket{\tilde p_i^{a,\kk}}$ refers to the Bloch sum, not to
any $\kk$-dependence of the underlying atomic shape. This distinction will be
important when we later discuss rotating projectors in $G$-space.

\subsection{Plane-wave representation}

For practical calculations (and for everything that follows), it is
convenient to expand pseudo-orbitals and Bloch projectors in plane waves:
\begin{align}
\tilde\psi_{n\kk}(\rr)
  &= \frac{1}{\sqrt{\Omega}} \sum_{\GG}
      \tilde C_{n\kk}(\GG)\, e^{i(\kk + \GG)\cdot\rr}, \\
\tilde p_i^{a,\kk}(\rr)
  &= \frac{1}{\sqrt{\Omega}} \sum_{\GG}
      \tilde P_i^{a,\kk}(\GG)\, e^{i(\kk + \GG)\cdot\rr},
\end{align}
where $\Omega$ is the cell volume. With our convention
\eqref{eq:bloch-projector}, an explicit calculation of the Fourier transform of
a single localized projector yields
\begin{equation}
\tilde P_i^{a,\kk}(\GG)
  \;=\; e^{i(\kk + \GG)\cdot\tautau_a}\,
        \tilde{\mathfrak p}_i^{a}(\kk + \GG),
\label{eq:projector-pw}
\end{equation}
where $\tilde{\mathfrak p}_i^{a}(\mathbf{q})$ is the species-only Fourier
transform of the atomic projector centered at the origin --- it has no
$\kk$-dependence and no site-position phase.

The projector coefficient is then a plane-wave inner product,
\begin{equation}
c_i^a(n,\kk)
  \;=\; \sum_{\GG} \tilde P_i^{a,\kk}(\GG)^*\,\tilde C_{n\kk}(\GG).
\label{eq:projector-pw-inner}
\end{equation}

\section{Real-space symmetry relation for the AE orbital}

Let $\kk' = R\kk$. Since the AE Hamiltonian commutes with $\Uhat(g)$ for
every $g$ in the space group, the AE orbital at $\kk'$ is, up to a phase
convention,
\begin{equation}
\ket{\psi_{n\kk'}} \;=\; \Uhat(g)\,\ket{\psi_{n\kk}}.
\label{eq:AE-symm}
\end{equation}
Applying $\Uhat(g)$ term by term to \eqref{eq:paw-bloch} gives, schematically,
\begin{equation}
\ket{\psi_{n\kk'}}
  \;=\; \Uhat(g)\,\ket{\tilde\psi_{n\kk}}
  + \sum_{a,i}\Uhat(g)\,
    \Bigl(\ket{\phi_i^{a,\kk}} - \ket{\tilde\phi_i^{a,\kk}}\Bigr)\,
    c_i^a(n,\kk).
\end{equation}
We must now compute how each ingredient transforms.

\subsection{Pseudo-orbital}

The pseudo-Hamiltonian carries the same symmetry, so
\begin{equation}
\ket{\tilde\psi_{n\kk'}} \;=\; \Uhat(g)\,\ket{\tilde\psi_{n\kk}},
\qquad
\tilde\psi_{n\kk'}(\rr) \;=\; \tilde\psi_{n\kk}\bigl(R^{-1}(\rr - \tvec)\bigr).
\label{eq:pseudo-symm}
\end{equation}

\subsection{Localized partial waves and projectors}

Apply $\Uhat(g)$ to a localized partial wave at $(a,\RR)$. Writing
$\rr_a^{\RR} \equiv \rr - \tautau_a - \RR$, the transformed function evaluates
the original at $g^{-1}\rr$:
\begin{align}
\bigl[\Uhat(g)\phi_i^{a,\RR}\bigr](\rr)
  &= \phi_i^{a,\RR}\!\bigl(R^{-1}(\rr - \tvec)\bigr) \nonumber\\
  &= R_{n_i\ell_i}^a\bigl(|R^{-1}(\rr - \tvec) - \tautau_a - \RR|\bigr)\,
     Y_{\ell_i m_i}\bigl(\widehat{R^{-1}(\rr - \tvec) - \tautau_a - \RR}\bigr).
\end{align}
The argument inside the radial function simplifies because $R \in O(3)$:
\begin{equation}
\bigl|R^{-1}(\rr - \tvec) - \tautau_a - \RR\bigr|
  \;=\; \bigl|\rr - \tvec - R\tautau_a - R\RR\bigr|
  \;=\; \bigl|\rr - (R\tautau_a + \tvec) - R\RR\bigr|.
\end{equation}

\paragraph{Site permutation.} The point $R\tautau_a + \tvec$ is, by the
definition of a space-group operation, equivalent to some basis position
$\tautau_b$ in the home cell up to a lattice translation. Define the
\emph{site permutation} $b = g(a)$ and the \emph{folding lattice vector}
$\TT_a$ via
\begin{equation}
R\,\tautau_a + \tvec \;=\; \tautau_{g(a)} + \TT_a, \qquad \TT_a \in \text{Bravais lattice}.
\label{eq:site-permutation}
\end{equation}
Then
\begin{equation}
\rr - (R\tautau_a + \tvec) - R\RR
  \;=\; \rr - \tautau_{g(a)} - (R\RR + \TT_a),
\end{equation}
so the rotated partial wave is centered at site $g(a)$ in cell
$\RR' = R\RR + \TT_a$:
\begin{equation}
\bigl[\Uhat(g)\phi_i^{a,\RR}\bigr](\rr)
  \;=\; R_{n_i\ell_i}^{g(a)}\bigl(|\rr - \tautau_{g(a)} - \RR'|\bigr)\,
        Y_{\ell_i m_i}\!\Bigl(R^{-1}\,\widehat{\rr - \tautau_{g(a)} - \RR'}\Bigr).
\label{eq:rotated-partial}
\end{equation}
The radial function attaches to species $g(a)$, which is the same species as
$a$ (otherwise $g$ would not be a space-group operation), so the radial part
is unchanged.

\paragraph{Wigner $D$-matrix on the angular part.} Real or complex spherical
harmonics transform under $R$ as
\begin{equation}
Y_{\ell m}(R^{-1}\hat{\mathbf n})
  \;=\; \sum_{m'} D^{(\ell)}_{m'm}(R)\,Y_{\ell m'}(\hat{\mathbf n}),
\label{eq:wigner-Y}
\end{equation}
where $D^{(\ell)}(R)$ is the $(2\ell+1)\times(2\ell+1)$ Wigner rotation matrix
in whichever basis (real or complex spherical harmonics) the code uses.
Substituting \eqref{eq:wigner-Y} into \eqref{eq:rotated-partial} gives
\begin{equation}
\Uhat(g)\,\ket{\phi_i^{a,\RR}}
  \;=\; \sum_{m'} D^{(\ell_i)}_{m' m_i}(R)\,
        \ket{\phi_{(n_i\ell_i m')}^{g(a),\,R\RR + \TT_a}}.
\label{eq:partial-localized-symm}
\end{equation}
The projector $\tilde p_i^{a,\RR}$ transforms identically (it has the same
angular structure):
\begin{equation}
\Uhat(g)\,\ket{\tilde p_i^{a,\RR}}
  \;=\; \sum_{m'} D^{(\ell_i)}_{m' m_i}(R)\,
        \ket{\tilde p_{(n_i\ell_i m')}^{g(a),\,R\RR + \TT_a}}.
\label{eq:projector-localized-symm}
\end{equation}

\subsection{Bloch sum after rotation}

Apply $\Uhat(g)$ to the Bloch-summed projector \eqref{eq:bloch-projector}:
\begin{align}
\Uhat(g)\,\ket{\tilde p_i^{a,\kk}}
  &= \sum_{\RR} e^{i\kk\cdot(\tautau_a + \RR)}\,\Uhat(g)\,\ket{\tilde p_i^{a,\RR}} \nonumber\\
  &= \sum_{\RR}\sum_{m'} e^{i\kk\cdot(\tautau_a + \RR)}\,
     D^{(\ell_i)}_{m' m_i}(R)\,
     \ket{\tilde p_{(n_i\ell_i m')}^{g(a),\,R\RR + \TT_a}}.
\label{eq:projector-bloch-1}
\end{align}
Change summation variable to $\RR' = R\RR + \TT_a$, equivalently
$\RR = R^{-1}(\RR' - \TT_a)$. Using
$\kk\cdot R^{-1}\mathbf{v} = (R\kk)\cdot\mathbf{v} = \kk'\cdot\mathbf{v}$ for
any vector $\mathbf{v}$, the phase becomes
\begin{equation}
\kk\cdot(\tautau_a + \RR)
  \;=\; \kk\cdot\tautau_a + \kk'\cdot(\RR' - \TT_a)
  \;=\; \kk'\cdot\RR' - \kk'\cdot\TT_a + \kk\cdot\tautau_a.
\end{equation}
Now we want to recognize the Bloch sum at $\kk'$ centered at site $g(a)$,
which is
\begin{equation}
\ket{\tilde p_j^{g(a),\kk'}}
  \;=\; \sum_{\RR'} e^{i\kk'\cdot(\tautau_{g(a)} + \RR')}\,\ket{\tilde p_j^{g(a),\RR'}}.
\end{equation}
Comparing, we need to factor out $e^{i\kk'\cdot\tautau_{g(a)}}$ and a leftover
phase. Using \eqref{eq:site-permutation}, $\tautau_{g(a)} = R\tautau_a + \tvec - \TT_a$,
so
\begin{equation}
\kk'\cdot\tautau_{g(a)} \;=\; \kk'\cdot(R\tautau_a) + \kk'\cdot\tvec - \kk'\cdot\TT_a
  \;=\; \kk\cdot\tautau_a + \kk'\cdot\tvec - \kk'\cdot\TT_a.
\end{equation}
Therefore
\begin{align}
\kk\cdot(\tautau_a + \RR)
  &= \kk'\cdot\RR' - \kk'\cdot\TT_a + \kk\cdot\tautau_a \nonumber\\
  &= \kk'\cdot\RR' + \kk'\cdot\tautau_{g(a)} - \kk'\cdot\tvec \nonumber\\
  &= \kk'\cdot(\tautau_{g(a)} + \RR') - \kk'\cdot\tvec.
\end{align}
Substituting into \eqref{eq:projector-bloch-1},
\begin{equation}
\boxed{\;\Uhat(g)\,\ket{\tilde p_i^{a,\kk}}
  \;=\; e^{-i\kk'\cdot\tvec}\,\sum_{m'} D^{(\ell_i)}_{m' m_i}(R)\,
        \ket{\tilde p_{(n_i\ell_i m')}^{g(a),\kk'}}.\;}
\label{eq:projector-bloch-symm}
\end{equation}
Identical statements hold for $\phi_i^{a,\kk}$ and $\tilde\phi_i^{a,\kk}$.
The fractional-translation phase $e^{-i\kk'\cdot\tvec}$ is global (band- and
site-independent); the Wigner matrix $D^{(\ell_i)}(R)$ mixes the $m$-channels;
and the site label is permuted from $a$ to $g(a)$. Note that, in this
convention, the lattice-folding phase $e^{-i\kk\cdot\TT_a}$ that appeared in
earlier intermediate expressions has been absorbed into the basis-position
phase of the Bloch projector at $\kk'$.

\subsection{Projector-coefficient transformation}

Insert \eqref{eq:projector-bloch-symm} and \eqref{eq:pseudo-symm} into the
projector coefficient at $\kk'$, using unitarity of $\Uhat(g)$:
\begin{align}
c_j^{b}(n,\kk')
  &= \langle \tilde p_j^{b,\kk'} \mid \tilde\psi_{n\kk'}\rangle
   = \langle \tilde p_j^{b,\kk'} \mid \Uhat(g)\,\tilde\psi_{n\kk}\rangle \nonumber\\
  &= \langle \Uhat(g)^\dagger\,\tilde p_j^{b,\kk'} \mid \tilde\psi_{n\kk}\rangle.
\end{align}
Apply \eqref{eq:projector-bloch-symm} \emph{backwards}: with $b = g(a)$ and
$\Uhat(g)^\dagger = \Uhat(g^{-1})$,
\begin{equation}
\Uhat(g)^\dagger\,\ket{\tilde p_j^{g(a),\kk'}}
  \;=\; e^{+i\kk\cdot R^{-1}\tvec}\,\sum_{m} \bigl[D^{(\ell_j)}(R)^\dagger\bigr]_{m m_j}\,
        \ket{\tilde p_{(n_j\ell_j m)}^{a,\kk}}.
\end{equation}
Take the conjugate of the inner product and rearrange. The cleaner way is to
\emph{forward}-substitute: from \eqref{eq:projector-bloch-symm},
\begin{equation}
\ket{\tilde p_j^{g(a),\kk'}}
  \;=\; e^{+i\kk'\cdot\tvec}\,\sum_{m}
        \bigl[D^{(\ell_j)}(R)^{-1}\bigr]_{m m_j}\,
        \Uhat(g)\,\ket{\tilde p_{(n_j\ell_j m)}^{a,\kk}},
\end{equation}
so that
\begin{align}
c_{(n_j\ell_j m_j)}^{g(a)}(n,\kk')
  &= e^{-i\kk'\cdot\tvec}\,
     \sum_m \bigl[D^{(\ell_j)}(R)\bigr]_{m_j m}^{*}\,
     \langle \Uhat(g)\,\tilde p_{(n_j\ell_j m)}^{a,\kk}\mid\Uhat(g)\,\tilde\psi_{n\kk}\rangle \nonumber\\
  &= e^{-i\kk'\cdot\tvec}\,
     \sum_m \bigl[D^{(\ell_j)}(R)^{*}\bigr]_{m_j m}\,
     c_{(n_j\ell_j m)}^{a}(n,\kk).
\end{align}
Using unitarity of $D^{(\ell)}$, $[D^{(\ell)}(R)^*]_{m_j m} = [D^{(\ell)}(R^{-1})]_{m m_j}^{T}
 = D^{(\ell)}(R)_{m m_j}^{*}$, we can write the result more transparently as
\begin{equation}
\boxed{\;
c_{(n\ell m')}^{g(a)}(n,\kk')
  \;=\; e^{-i\kk'\cdot\tvec}\,
        \sum_m D^{(\ell)}_{m' m}(R)\,c_{(n\ell m)}^{a}(n,\kk).
\;}
\label{eq:coef-symm}
\end{equation}
The radial-channel index $n\ell$ is preserved; only $m$ rotates and the
site label is permuted.

If a different gauge is chosen for the Bloch projector --- for instance, the
\emph{lattice-only} convention without the $e^{i\kk\cdot\tautau_a}$ factor ---
the formula \eqref{eq:coef-symm} acquires an additional site-folding phase
$e^{-i\kk\cdot\TT_a}$. The boxed form here corresponds to convention
\eqref{eq:bloch-projector}.

\section{Generating $R\kk$ from $\kk$ in $G$-space}

\subsection{Plane-wave transformation rule}

We now derive the $G$-space rule used in practice. Consider a generic Bloch
function $f_\kk(\rr) = \sum_\GG F_\kk(\GG)\, e^{i(\kk+\GG)\cdot\rr}$. Apply
$\Uhat(g)$:
\begin{align}
\bigl[\Uhat(g) f_\kk\bigr](\rr)
  &= f_\kk\bigl(R^{-1}(\rr - \tvec)\bigr) \nonumber\\
  &= \sum_\GG F_\kk(\GG)\, e^{i(\kk + \GG)\cdot R^{-1}(\rr - \tvec)} \nonumber\\
  &= \sum_\GG F_\kk(\GG)\, e^{-i(R\kk + R\GG)\cdot\tvec}\,
                       e^{i(R\kk + R\GG)\cdot\rr} \nonumber\\
  &= \sum_{\GG'} F_\kk(R^{-1}\GG')\, e^{-i(\kk' + \GG')\cdot\tvec}\,
                       e^{i(\kk' + \GG')\cdot\rr},
\end{align}
where $\GG' = R\GG$ and we used $\kk' = R\kk$. Reading off the coefficients,
\begin{equation}
\boxed{\;
\bigl[\Uhat(g) f_\kk\bigr]_{\kk'}\!(\GG')
  \;=\; e^{-i(\kk' + \GG')\cdot\tvec}\,F_\kk\bigl(R^{-1}\GG'\bigr).
\;}
\label{eq:G-space-rotation}
\end{equation}
Equivalently, $F_{\kk'}(R\GG) = e^{-i(\kk' + R\GG)\cdot\tvec}\,F_\kk(\GG)$.
Three observations:
\begin{enumerate}[leftmargin=*]
\item The map $\GG \mapsto R\GG$ permutes vectors of the $\kk$-sphere onto
      vectors of the $\kk'$-sphere; both spheres have the same radius because
      $|R(\kk+\GG)| = |\kk + \GG|$.
\item The phase $e^{-i(\kk'+\GG')\cdot\tvec}$ vanishes for symmorphic operations
      ($\tvec = 0$).
\item If $R\kk$ differs from the canonical $\kk'$ in the BZ by a reciprocal
      lattice vector $\GG_0$, $R\kk = \kk' + \GG_0$, then one uses
      $\GG' = R\GG - \GG_0$ so that $\kk' + \GG' = R(\kk + \GG)$. An overall
      band phase $e^{-i\GG_0\cdot\tvec}$ may then be absorbed by convention.
\end{enumerate}

\subsection{Application to pseudo-orbital and projector}

The pseudo-orbital coefficients transform via \eqref{eq:G-space-rotation}:
\begin{equation}
\tilde C_{n\kk'}(R\GG) \;=\; e^{-i(\kk'+R\GG)\cdot\tvec}\,\tilde C_{n\kk}(\GG).
\label{eq:pseudo-PW-symm}
\end{equation}
Similarly for the Bloch-summed projector --- but with one extra wrinkle.
Because $\Uhat(g)$ acts on a localized projector by both rotating its angular
shape \emph{and} relocating its center to site $g(a)$, the $G$-space rotation
rule connects the projector at site $a$ at $\kk$ to the rotated function,
which is a Wigner-mixed combination of native projectors at site $g(a)$ at
$\kk'$:
\begin{equation}
\bigl[\Uhat(g)\tilde p_i^{a,\kk}\bigr]_{\kk'}\!(\GG')
  \;=\; e^{-i(\kk'+\GG')\cdot\tvec}\,\tilde P_i^{a,\kk}\bigl(R^{-1}\GG'\bigr).
\label{eq:projector-PW-rotation}
\end{equation}
Equivalently, comparing with \eqref{eq:projector-bloch-symm} in the plane-wave
representation, the right-hand side equals
\begin{equation}
e^{-i\kk'\cdot\tvec}\,\sum_{m'} D^{(\ell_i)}_{m' m_i}(R)\,
      \tilde P^{g(a),\kk'}_{(n_i\ell_i m')}(\GG'),
\end{equation}
which gives a useful identity satisfied by the species-only function
$\tilde{\mathfrak p}^a$:
\begin{equation}
\tilde{\mathfrak p}^a_i(R^{-1}\mathbf{q})
  \;=\; \sum_{m'} D^{(\ell_i)}_{m' m_i}(R)\,
        \tilde{\mathfrak p}^a_{(n_i\ell_i m')}(\mathbf{q}),
\end{equation}
which is just the statement that the angular part is a single spherical
harmonic $Y_{\ell_i m_i}$.

\subsection{View 1 vs.~View 2 of the on-site basis}

Two consistent conventions are possible at $\kk'$:
\begin{description}[leftmargin=2em]
\item[\textbf{View 1 (rotate everything together).}]
  Define the projector at site $g(a)$ at $\kk'$ as the
  $G$-rotated image of the projector at site $a$ at $\kk$, i.e.\
  use \eqref{eq:projector-PW-rotation} as the \emph{definition}.
  Under this convention,
  \begin{equation}
  c_i^{g(a)}(n,\kk') \;=\; c_i^{a}(n,\kk),
  \label{eq:view1-equality}
  \end{equation}
  with no Wigner matrix and no fractional-translation phase: the inner
  product is invariant under the same unitary applied to bra and ket. This is
  the natural representation produced by a $G$-space-only algorithm, and we
  adopt it from Section~\ref{sec:multi-symm} onward.
\item[\textbf{View 2 (canonical $\ell m$ basis at every $\kk$).}]
  Insist that the projector at site $b$ in channel $(n\ell m)$ is always the
  Bloch sum of the same atomic shape with definite $(\ell,m)$ in the global
  frame. Then the coefficients transform via
  \eqref{eq:coef-symm} with explicit Wigner matrices.
\end{description}
The two are related by a unitary basis change on the on-site space; physical
quantities (densities, energies, AE orbitals) are identical.

\subsection{Practical algorithm (View 1)}

Given $\kk$-data $\{\tilde C_{n\kk}(\GG),\,\tilde P_i^{a,\kk}(\GG)\}$ and a
symmetry $g = \{R\mid\tvec\}$:
\begin{enumerate}[leftmargin=*]
\item Build the $G$-permutation $\GG\mapsto R\GG$ once. (If $R\kk = \kk' +
      \GG_0$, fold by $\GG_0$.)
\item For each band, apply \eqref{eq:pseudo-PW-symm} to obtain
      $\tilde C_{n\kk'}(\GG')$.
\item For each site $a$ and channel $i$, apply
      \eqref{eq:projector-PW-rotation} to obtain
      $\bigl[\Uhat(g)\tilde p_i^{a,\kk}\bigr]_{\kk'}\!(\GG')$, and label this
      as the projector at site $g(a)$ at $\kk'$ (View 1).
\item Compute the projector coefficient at $\kk'$ by the ordinary
      $G$-space inner product
      \begin{equation}
      c_i^{g(a)}(n,\kk')
        \;=\; \sum_{\GG'} \bigl[\Uhat(g)\tilde p_i^{a,\kk}\bigr]_{\kk'}\!(\GG')^*\,
              \tilde C_{n\kk'}(\GG').
      \end{equation}
      By construction this equals $c_i^a(n,\kk)$.
\end{enumerate}
The recipe never requires explicit Wigner matrices. They are still acting,
implicitly, in how the angular content of the projector reshuffles when
$\GG\to R\GG$, but the bookkeeping is hidden inside the $G$-space rotation.

\section{Two operations sharing a $\kibz$: little-group representation}
\label{sec:multi-symm}

Suppose we have computed the orbitals at $\kibz$ and generate orbitals at two
star-related $\kk$-points,
\begin{equation}
\kk_1 \;=\; R_1\,\kibz, \qquad \kk_2 \;=\; R_2\,\kibz,
\end{equation}
by the View-1 procedure above. Question: what is the expansion of an orbital
at $\kk_1$ in the basis of orbitals at $\kk_2$?

\subsection{Crystal-momentum constraint}

For Bloch states, $\langle\psi_{m\kk_2}\mid\psi_{n\kk_1}\rangle = 0$ unless
$\kk_1 \equiv \kk_2 \pmod{\text{reciprocal lattice}}$. Otherwise, $S_{mn} =
0$ and the expansion is empty. The non-trivial case is therefore
\begin{equation}
\kk_1 \equiv \kk_2 \pmod{\text{recip.~lat.}}
  \;\Longleftrightarrow\;
  R_2^{-1} R_1 \in G_{\kk_{\mathrm{ibz}}},
\end{equation}
where $G_{\kk_{\mathrm{ibz}}}$ is the little group of $\kibz$ (point-group operations that
fix $\kibz$ modulo reciprocal lattice vectors). Define
\begin{equation}
g_0 \;\equiv\; g_2^{-1}\,g_1
   \;=\; \{R_2^{-1}R_1 \mid R_2^{-1}(\tt_1 - \tt_2)\}
   \;=\; \{R_0 \mid \tt_0\},
\end{equation}
where $R_0 = R_2^{-1}R_1$ and $\tt_0 = R_2^{-1}(\tt_1 - \tt_2)$. (We assume
the symmorphic case $\tt_1 = \tt_2 = 0$ if the reader prefers; non-trivial
$\tvec$'s only add explicit phases via \eqref{eq:G-space-rotation}.)

\subsection{Reduction to the little-group matrix at $\kibz$}

By construction in View 1,
\begin{equation}
\ket{\psi_{n\kk_1}} \;=\; \Uhat(g_1)\,\ket{\psi_{n\kibz}},
\qquad
\ket{\psi_{m\kk_2}} \;=\; \Uhat(g_2)\,\ket{\psi_{m\kibz}}.
\end{equation}
Therefore the overlap factorizes via the unitarity of $\Uhat$:
\begin{align}
S_{mn}
  &\;\equiv\; \langle\psi_{m\kk_2}\mid\psi_{n\kk_1}\rangle \nonumber\\
  &= \langle \Uhat(g_2)\,\psi_{m\kibz}\mid \Uhat(g_1)\,\psi_{n\kibz}\rangle \nonumber\\
  &= \langle \psi_{m\kibz}\mid \Uhat(g_2)^\dagger\,\Uhat(g_1)\,\psi_{n\kibz}\rangle \nonumber\\
  &= \langle \psi_{m\kibz}\mid \Uhat(g_0)\,\psi_{n\kibz}\rangle.
\end{align}
That is,
\begin{equation}
\boxed{\;
S_{mn} \;=\; D_{mn}(g_0),
\qquad
D_{mn}(g) \;\equiv\; \langle\psi_{m\kibz}\mid\Uhat(g)\,\psi_{n\kibz}\rangle,
\quad g\in G_{\kk_{\mathrm{ibz}}}.
\;}
\label{eq:S-as-D}
\end{equation}
This is the matrix of the little-group operation $g_0$ in the band basis at
$\kibz$. In particular, $D(g)$ is a unitary representation (projective in the
nonsymmorphic case) of $G_{\kk_{\mathrm{ibz}}}$ on the Bloch states at $\kibz$.

\subsection{Computing $D(g_0)$ in the PAW $G$-space representation}

Using the full PAW overlap (with on-site augmentation),
\begin{equation}
\langle\psi_{m\kibz}\mid\psi'_{n\kibz}\rangle
  \;=\; \langle\tilde\psi_{m\kibz}\mid\tilde\psi'_{n\kibz}\rangle
   + \sum_{a,ij} c_i^{a}(m,\kibz)^*\,Q^a_{ij}\,c_j^{a,\,\prime}(n,\kibz),
\end{equation}
where the augmentation matrix
\begin{equation}
Q^a_{ij} \;=\; \langle\phi_i^{a}\mid\phi_j^{a}\rangle - \langle\tilde\phi_i^{a}\mid\tilde\phi_j^{a}\rangle
\end{equation}
is a fixed property of the PAW dataset, and we set $\ket{\psi'_{n\kibz}} =
\Uhat(g_0)\,\ket{\psi_{n\kibz}}$. The primed coefficients are
\begin{equation}
c_j^{a,\,\prime}(n,\kibz)
  \;=\; \langle\tilde p_j^{a,\kibz}\mid\Uhat(g_0)\,\tilde\psi_{n\kibz}\rangle.
\end{equation}

In $G$-space, with $\tilde C^{(g_0)}_{n}(\GG')$ denoting the coefficients of
the rotated pseudo-orbital,
\begin{equation}
\tilde C^{(g_0)}_{n}(R_0\GG)
  \;=\; e^{-i(\kibz + R_0\GG)\cdot\tt_0}\,\tilde C_{n\kibz}(\GG),
\end{equation}
we have
\begin{align}
\langle\tilde\psi_{m\kibz}\mid\Uhat(g_0)\,\tilde\psi_{n\kibz}\rangle
  &= \sum_{\GG} \tilde C_{m\kibz}(\GG)^*\,\tilde C^{(g_0)}_{n}(\GG), \\
c_j^{a,\,\prime}(n,\kibz)
  &= \sum_{\GG} \tilde P_j^{a,\kibz}(\GG)^*\,\tilde C^{(g_0)}_{n}(\GG).
\end{align}
Putting it together:
\begin{equation}
\boxed{\;
\begin{aligned}
S_{mn} \;=\; D_{mn}(g_0)
  &= \sum_{\GG} \tilde C_{m\kibz}(\GG)^*\,\tilde C^{(g_0)}_{n}(\GG) \\
  &\quad + \sum_{a,ij} c_i^a(m,\kibz)^*\,Q^a_{ij}\,c_j^{a,\,\prime}(n,\kibz),
\end{aligned}
\;}
\label{eq:S-final}
\end{equation}
where the rotated pseudo-orbital is built once via the $G$-space rotation
rule \eqref{eq:G-space-rotation}, and $c^{\prime}$ is then obtained as a
single inner product against the native projectors at $\kibz$.

\subsection{Special cases}

\paragraph{Non-degenerate band.} If band $n$ is non-degenerate at $\kibz$ and
sits in a one-dimensional irrep of $G_{\kk_{\mathrm{ibz}}}$, then
$\Uhat(g_0)\ket{\psi_{n\kibz}} = \lambda_n(g_0)\ket{\psi_{n\kibz}}$ with
$|\lambda_n(g_0)|=1$, and
\begin{equation}
\ket{\psi_{n\kk_1}} \;=\; \lambda_n(g_0)\,\ket{\psi_{n\kk_2}}.
\end{equation}

\paragraph{Degenerate multiplet.} If band $n$ lives in a $d$-dimensional
multiplet $\{\ket{\psi_{m\kibz}}\}_{m=1}^{d}$ that carries an irrep of
$G_{\kk_{\mathrm{ibz}}}$, then $D(g_0)$ restricted to the multiplet is the unitary
representation matrix of $g_0$ in that irrep. Diagonalizing $D(g)$ across
$g\in G_{\kk_{\mathrm{ibz}}}$ yields symmetry-adapted bands.

\paragraph{Trivial coincidence $g_1 = g_2$.} If $R_1 = R_2$, then $g_0 = e$,
$D(e)_{mn} = \delta_{mn}$, and $S_{mn} = \delta_{mn}$, so
$\ket{\psi_{n\kk_1}} = \ket{\psi_{n\kk_2}}$ identically.

\subsection{Composition law and consistency check}

Since $\Uhat$ is a representation of the space group (projective for
nonsymmorphic groups due to the fractional-translation phase), the matrices
$D(g)$ satisfy
\begin{equation}
D(g)\,D(g') \;=\; \omega(g,g')\,D(g\,g'),
\end{equation}
with $|\omega(g,g')|=1$. In the symmorphic case $\omega \equiv 1$. This
composition law is a stringent numerical test of one's $G$-space rotation
conventions: any error in the $\tvec$-phase or $\GG_0$-folding bookkeeping shows
up immediately as a violation of $D(g)D(g') = \pm D(gg')$ at machine
precision.

\section{Concluding remarks: spin and time reversal}

\paragraph{Spin--orbit coupling.} For spinor wavefunctions, replace the
Wigner $D^{(\ell)}$ matrices by spin--orbit-coupled $D^{(j)}$ matrices acting
on the appropriate $\ell s\to j$ basis. The two-component spinor part of the
plane-wave coefficients is rotated by the SU(2) matrix $u(R) =
\exp(-i\boldsymbol\sigma\cdot\hat{\mathbf{n}}\,\theta/2)$ associated with $R$.
The structure of \eqref{eq:G-space-rotation} and \eqref{eq:S-as-D} is
unchanged.

\paragraph{Time reversal.} The antiunitary operator $\Theta$ acts as
$[\Theta\psi_{n\kk}](\rr) = \psi_{n\kk}(\rr)^*$ in the spinless case, sending
$\kk \to -\kk$, complex-conjugating $\tilde C_{n\kk}(\GG)$, and conjugating
the projector coefficients. With SOC, $\Theta = i\sigma_y K$. Combining $g$
with $\Theta$ extends the available star of an irreducible $\kk$-point and is
treated by the same machinery, with antiunitary care for the inner-product
ordering.

\paragraph{Why this works in PAW.} The decisive structural fact is that, in
PAW, projectors and partial waves transform identically under space-group
operations: they share the same angular momentum representation and the same
site-permutation rule. Hence the augmentation contribution to the AE overlap
\emph{automatically} respects the symmetry, and the $G$-space rotation alone
is sufficient to relate orbitals at symmetry-equivalent $\kk$-points without
ever decomposing into Wigner $D$-matrices.

\end{document}
